%------------------------------------------------------------------------------%
\section{Volume pelo método das seções transversais}
\label{sec:Volume_secoes_transversais}

%------------------------------------------------------------------------------%
\begin{figure}
  \centering
  \includegraphics[width=0.5\linewidth]{volume-secoes-01}
  \caption{Volume por seções transversais -- 01}
  \label{fig:volume_secoes_01}
\end{figure}
%------------------------------------------------------------------------------%

%------------------------------------------------------------------------------%
\begin{figure}
  \centering
  \includegraphics[width=0.5\linewidth]{volume-secoes-02}
  \caption{Volume por seções transversais -- 02}
  \label{fig:volume_secoes_02}
\end{figure}
%------------------------------------------------------------------------------%

%------------------------------------------------------------------------------%
\begin{figure}
  \centering
  \includegraphics[width=0.5\linewidth]{volume-secoes-03}
  \caption{Volume por seções transversais -- 03}
  \label{fig:volume_secoes_03}
\end{figure}
%------------------------------------------------------------------------------%

%------------------------------------------------------------------------------%
\begin{figure}
  \centering
  \includegraphics[width=0.5\linewidth]{volume-secoes-04}
  \caption{Volume por seções transversais -- 04}
  \label{fig:volume_secoes_04}
\end{figure}
%------------------------------------------------------------------------------%

%------------------------------------------------------------------------------%
\begin{figure}
  \centering
  \includegraphics[width=0.5\linewidth]{volume-secoes-05}
  \caption{Volume por seções transversais -- 05}
  \label{fig:volume_secoes_05}
\end{figure}
%------------------------------------------------------------------------------%

%------------------------------------------------------------------------------%
\begin{figure}
  \centering
  \includegraphics[width=0.5\linewidth]{volume-secoes-06}
  \caption{Volume por seções transversais -- 06}
  \label{fig:volume_secoes_06}
\end{figure}
%------------------------------------------------------------------------------%

%------------------------------------------------------------------------------%
\begin{figure}
  \centering
  \includegraphics[width=0.5\linewidth]{volume-secoes-07}
  \caption{Volume por seções transversais -- 07}
  \label{fig:volume_secoes_01}
\end{figure}
%------------------------------------------------------------------------------%

%------------------------------------------------------------------------------%
\begin{figure}
  \centering
  \includegraphics[width=0.5\linewidth]{volume-secoes-08}
  \caption{Volume por seções transversais -- 08}
  \label{fig:volume_secoes_08}
\end{figure}
%------------------------------------------------------------------------------%

%------------------------------------------------------------------------------%
\begin{figure}
  \centering
  \includegraphics[width=0.5\linewidth]{volume-secoes-09}
  \caption{Volume por seções transversais -- 09}
  \label{fig:volume_secoes_09}
\end{figure}
%------------------------------------------------------------------------------%

%------------------------------------------------------------------------------%
\begin{figure}
  \centering
  \includegraphics[width=0.5\linewidth]{volume-secoes-10}
  \caption{Volume por seções transversais -- 10}
  \label{fig:volume_secoes_10}
\end{figure}
%------------------------------------------------------------------------------%

%------------------------------------------------------------------------------%
\begin{figure}
  \centering
  \includegraphics[width=0.5\linewidth]{volume-secoes-11}
  \caption{Volume por seções transversais -- 11}
  \label{fig:volume_secoes_11}
\end{figure}
%------------------------------------------------------------------------------%

%------------------------------------------------------------------------------%
\begin{figure}
  \centering
  \includegraphics[width=0.5\linewidth]{volume-secoes-12}
  \caption{Volume por seções transversais -- 12}
  \label{fig:volume_secoes_12}
\end{figure}
%------------------------------------------------------------------------------%

%------------------------------------------------------------------------------%
\begin{figure}
  \centering
  \includegraphics[width=0.5\linewidth]{volume-secoes-13}
  \caption{Volume por seções transversais -- 13}
  \label{fig:volume_secoes_13}
\end{figure}
%------------------------------------------------------------------------------%

%------------------------------------------------------------------------------%
\begin{figure}
  \centering
  \includegraphics[width=0.5\linewidth]{volume-secoes-14}
  \caption{Volume por seções transversais -- 14}
  \label{fig:volume_secoes_14}
\end{figure}
%------------------------------------------------------------------------------%

%------------------------------------------------------------------------------%
\begin{figure}
  \centering
  \includegraphics[width=0.5\linewidth]{volume-secoes-15}
  \caption{Volume por seções transversais -- 15}
  \label{fig:volume_secoes_15}
\end{figure}
%------------------------------------------------------------------------------%

%------------------------------------------------------------------------------%
\begin{figure}
  \centering
  \includegraphics[width=0.5\linewidth]{volume-secoes-16}
  \caption{Volume por seções transversais -- 16}
  \label{fig:volume_secoes_16}
\end{figure}
%------------------------------------------------------------------------------%

%------------------------------------------------------------------------------%
\begin{figure}
  \centering
  \includegraphics[width=0.5\linewidth]{volume-secoes-17}
  \caption{Volume por seções transversais -- 17}
  \label{fig:volume_secoes_17}
\end{figure}
%------------------------------------------------------------------------------%

Nesta seção, iremos explorar uma aplicação importante das integrais no
cálculo de volumes de sólidos específicos. Em particular, abordaremos o
cálculo de volumes para cilindros e sólidos de rotação.

%------------------------------------------------------------------------------%
\begin{definicao}{Volume por Seções transversais}{teo:volsecoes}
  Seja $S$ um sólido que está definido para $a \leq x \leq b$. Se a área da
  secção transversal de $S$ no plano $P_x$ , passando por $x$ e
  perpendicular ao eixo $x$, é $A(x)$, em que A é uma função
  contínua, então o volume de $S$ é:
  \[
    V = \lim_{n \to \infty} \sum_{i=1}^{n} A(x_i) \Delta x_i
      = \int_a^b A(x)dx
  \]
  Essa definição é válida sempre que $A(x)$ for integrável, em particular, quando for contínua.
\end{definicao}
%------------------------------------------------------------------------------%

Para o cálculo do volume do sólido, siga os seguintes passos:
\begin{enumerate}
  \item Faça o esboço do sólido e uma seção transversal típica.
  \item Encontre uma expressão para $A(x)$, a área dessa seção transversal típica.
  \item Determine os limites de integração.
  \item Integre $A(x)$ para determinar o volume.
\end{enumerate}

%------------------------------------------------------------------------------%
\begin{example}
  Uma pirâmide de base quadrada com lado $3$ metros, tem $3$ metros de altura. A seção transversal da pirâmide, perpendicular a altura e a $x$ metros abaixo do vértice, é um quadrado com $x$ metros de lado. Determine o volume dessa pirâmide.
  \solution
  \[
    V
    \;=\; \int_0^3 A(x) dx
    \;=\; \int_0^3 x^2 dx
    \;\stackTFCD\; \left(\frac{x^3}{3}\right) \evalat{0}{3}
    \;=\; \frac{27}{3}
    \;=\; 9 \,\si{\meter^3}
  \]
\end{example}
%------------------------------------------------------------------------------%

%------------------------------------------------------------------------------%
\begin{example}
  Usando integrais, encontre o volume $V$ de uma esfera de raio $r$,
  mostrando que \(V= \frac{4}{3}\pi r^3\).
  \solution
  \begin{align*}
    V & = \int_{-r}^{r}A(x)dx                                            \\
      & = \int_{-r}^{r}\pi (r^2-x^2)dx                                   \\
      & = 2 \pi\int_{0}^{r} (r^2-x^2)dx                                  \\
      & \stackTFCD 2 \pi \left( r^2 x-\frac{x^3}{3}\right) \evalat{0}{r} \\
      & =2 \pi \left( r^3-\frac{r^3}{3}\right)                           \\
      & = \frac{4}{3} \pi r^3.
  \end{align*}
\end{example}
%------------------------------------------------------------------------------%

%------------------------------------------------------------------------------%
\begin{definicao}{Sólido de revolução}{def:solido-revolucao}
  Sólido de revolução é o sólido obtido pela rotação (ou revolução)
  de uma região plana em torno de um eixo. Nesse caso, a seção transversal
  será um disco de raio $R(x)$, que terá área
  \(
    A(x)= \pi R(x)^2
  \)
  ou um anel circular de raio externo $R(x)$ e raio interno $r(x)$, cuja área é
  \(
    \pi \left( R(x)^2 - r(x)^2 \right)
  \).
\end{definicao}
%------------------------------------------------------------------------------%

%------------------------------------------------------------------------------%
\begin{example}
  Encontre o volume do sólido obtido pela
  rotação em torno do eixo $x$ da região sob a curva
  $f(x)= \sqrt{x}$ de $0$ até $1$.
  \solution
  A seção transversal é um círculo de raio $R(x) = y = \sqrt{x}$ cuja área é dada por
  \[
    A(x)
    = \pi R(x)^2
    = \pi \left(\sqrt{x}\right)^2
    = \pi x \,.
  \]
  Além disso, o sólido está entre $0$ e $1$. Portanto
  \[
    V
    = \int_0^1 \pi x dx                                                                     \\
    \stackTFCD \pi \left( \frac{x^2}{2}\right) \evalat{0}{1}]
    = \frac{\pi}{2}.
  \]
\end{example}
%------------------------------------------------------------------------------%

Ao invés de rotacionarmos uma seção plana em torno do eixo $x$, podemos fazer
a rotação em torno do eixo $y$. Assim, teremos como região plana uma seção
circular de raio $x=R(y)$, $c \leq y \leq d$ e usamos o mesmo método trocando
$x$ por $y$. Ou seja, nesse caso, o volume do sólido obtido é dado por
\[
  V
  = \int_c^d A(y)dy
  = \int_c^d \pi R(y)^2 dy
\]
O exemplo a seguir ilustra essa situação.

%------------------------------------------------------------------------------%
\begin{example}
  Encontre o volume do sólido obtido pela
  rotação da região limitada por $y = x^3$, $y = 8$, e $x = 0$
  em torno do eixo $y$.
  \solution
  Observe que dessa vez, como a região é girada em torno do eixo $y$, faz
  sentido fatiar o sólido perpendicularmente ao eixo $y$, obtendo como
  seção transversal é um círculo de raio $x$, onde $x= \sqrt[3]{y}$.
  Assim, a área dessa seção transversal é dada por
  \[
    A(y)
    = \pi R(y)^2
    = \pi \left( \sqrt[3]{y} \right)^2
    = \pi y^{\sfrac{2}{3}}
  \]
  Ademais, uma vez que sólido está entre $y = 0$ e $y = 8$, seu
  volume é dado por:
  \[
    V
    = \int_0^8 \pi y^{\sfrac{2}{3}}dy
    \stackTFCD \pi \left( \frac{3}{5} y^{\sfrac{5}{3}}\right) \evalat{0}{8}
    = \pi \frac{3}{5} 8^{\sfrac{5}{3}}
    = \frac{96 \pi}{5}.
  \]
\end{example}
%------------------------------------------------------------------------------%

Em síntese, temos que:
\begin{itemize}
  \item Se o sólido é obtido pela rotação em torno do eixo $x$ de uma
  região plana cuja seção transversal é um círculo, o raio desse círculo
  é $y= f(x)$ e o volume é dado por uma integral na variável $x$.
  \item Por outro lado, se o sólido é obtido pela rotação em torno do
  eixo $y$ de uma região plana cuja seção transversal é um círculo, o
  raio desse círculo é $x= f(y)$ e o volume é dado por uma integral
  na variável $y$.
\end{itemize}

%------------------------------------------------------------------------------%
\begin{example}
  A região $\mathcal{R}$, limitada pelas curvas delimitada pelas curvas
  $y = x$ e $y = x^2$ é girada em torno do eixo x. Encontre o volume do
  sólido resultante.
  \solution
  A primeira observação importante é que as curvas $y = x$ e $y = x^2$
  se intersectam nos pontos $(0, 0)$ e $(1, 1)$. Esses valores são obtidos
  resolvendo a equação $x = x^2$.

  Nesse contexto, a seção transversal não é um círculo, mas sim um anel
  com raio interno $r(x) = x^2$ e raio externo $R(x) = x$.
  Dessa forma, a área da seção transversal é dada por
  \[
    A(x)
    = \pi \left( x^2 - \left(x^2\right)^2 \right)
    = \pi \left( x^2 - x^4 \right).
  \]
  Portanto:
  \[
    V
    = \int_0^1 \pi (x^2-x^4) dx
    \stackTFCD \pi \left( \frac{x^3}{3}-\frac{x^5}{5} \right) \evalat{0}{1}
    = \frac{1}{3} - \frac{1}{5}
    = \frac{2\pi}{15} .
    \]
\end{example}
%------------------------------------------------------------------------------%

Em vez de realizar a rotação da região plana em torno dos eixos coordenados,
podemos ter o interesse de calcular o volume de um sólido $S$ gerado pela
rotação de uma região plana em torno de uma reta paralela a esses eixos,
como ilustrado no exemplo a seguir.

\begin{example}
  Encontre o volume do sólido obtido pela rotação da região do
  exemplo anterior:
  \begin{enumerate}[label=\alph*)]
    \item Em torno da reta $x=-1$.
    \item Em torno da reta $y=2$
  \end{enumerate}
  \solution
  \begin{enumerate}[label=\alph*)]
    \item
    Devemos notar que, nessa situação, a seção transversal é um anel,
    mas dessa vez o raio interno $2-x$ e raio externo $2-x^2$.
    Assim, a área da seção transversal é dada por
    \[
      A(x)
      = \pi \left( 2 - x^2 \right)^2 - \pi (2-x)^2
      = \pi \left( x^4-5x^2+4x \right)
    \]
    Assim, o volume do sólido resultante é
    \begin{align*}
      V
      & = \int_0^1 \pi \left( x^4 - 5x^2 + 4x \right) dx \\
      & \stackTFCD \pi \left(
            \frac{x^5}{5} - 5 \frac{x^3}{3} + 4 \frac{x^2}{2}
          \right) \evalat{0}{1} \\
      & = \pi \left(\frac{1}{5} - \frac{5}{3} + \frac{4}{2}\right) \\
      & = \frac{8\pi}{15}
    \end{align*}

    \item
    Aqui, a seção transversal é um anel, com raio interno $1+y$ e
    raio externo $1+ \sqrt{y}$.
    Assim, a área da seção transversal é dada por
    \[
      A(y)
      = \pi \left( 1+\sqrt{y} \right)^2 - \pi (1+y)^2
      = \pi \left( 2\sqrt{y}-y-y^2 \right)
    \]
    Portanto, o volume pedido é dado por
    \begin{align*}
      V
      & = \int_0^1 \pi \left(2\sqrt{y}-y-y^2\right) dy \\
      & \stackTFCD \pi \left(
          \frac{4y^{\sfrac{3}{2}}}{3}-  \frac{y^2}{2}- \frac{y^3}{3}
        \right) \evalat{0}{1} \\
      & = \pi \left( \frac{4}{3} - \frac{1}{2} - \frac{1}{3}\right) \\
      & = \frac{\pi}{2}.
    \end{align*}
  \end{enumerate}
\end{example}
%------------------------------------------------------------------------------%

%------------------------------------------------------------------------------%
